% ====================================================
% SESSION 2 -- Questions, and where the data comes from
% Body file, shared by questions.tex and questions_notes.tex
% ====================================================

% ----------------------------------------------------
\begin{frame}[plain]
  \titlepage
\end{frame}
% ----------------------------------------------------
\note{}

% ----------------------------------------------------
\begin{frame}
\frametitle{Where we are}
\centering

\vspace{5pt}
\small
\begin{enumerate}
  \item \asher{Why design at all?} \hfill \asher{done}
  \vspace{6pt}
  \item \textbf{What is the question, and where would the data come from?} \hfill \textbf{today}
  \vspace{6pt}
  \item \asher{Do the numbers mean what you think they mean?}
  \vspace{6pt}
  \item \asher{What comparisons can we use to establish causal claims?}
  \vspace{6pt}
  \item \asher{What variation can you exploit, and how far does the answer travel?}
\end{enumerate}

\end{frame}
% ----------------------------------------------------
\note{One slide, ten seconds. It is what makes five lectures feel like one argument rather than a list of topics.

Today is two halves that look unrelated and are not: how to get to a question worth asking, and where the data that might answer it actually came from. The second half is the one nobody teaches, and it is where computational projects break.}

% ====================================================
\section{From topic to research question}
% ====================================================

% ----------------------------------------------------
\begin{frame}
\frametitle{Research and RQs}
\centering

\begin{itemize}
  \item Research $=$ answering questions
  \item<2-> Re-cap on research types:
  \begin{itemize}
    \item<2->[1.] \textbf{Normative} vs. \textbf{positive} research
    \item<3->[2.] Positive: \textbf{theoretical} vs. \textbf{empirical}
    \item<4->[3.] Empirical: \textbf{describe}, \textbf{explain}, \textbf{predict}
  \end{itemize}
\end{itemize}

\end{frame}
% ----------------------------------------------------
\note{Two minutes of placing. The third line is last week's triad, so this is a recap rather than new material.

The normative/positive line is worth thirty seconds: ``should Spain raise the minimum wage'' is not a research question in this sense, however important. ``What happened to employment when it was raised'' is. Students conflate the two constantly, and in the workshop it shows up as a design that cannot answer the thing they actually care about.}

% ----------------------------------------------------
\begin{frame}
\frametitle{Example on types of questions and evidence}
\centering

\begin{minipage}{.49\textwidth}\centering
  \begin{itemize}
    \item You are doing the analytics behind a campaign for Mahou
    \item[]
    \item Think of questions that are:
    \begin{itemize}
      \item Descriptive
      \item Explanatory
      \item Predictive
    \end{itemize}
    \item[]
    \item Which ones could you answer with the data a brewer actually has?
  \end{itemize}
\end{minipage}\hfill
\begin{minipage}{.49\textwidth}\centering

\includegraphics[width = 0.95\textwidth]{../img/mahou}
\end{minipage}

\end{frame}
% ----------------------------------------------------
\note{Pairs, four minutes, then take six or seven answers and sort them on the board into the three columns. It is the fastest way to find out whether the triad actually stuck from last week.

What usually comes back:
\begin{itemize}
  \item \textbf{Describe}: who drinks it, where, when, how has that changed?
  \item \textbf{Explain}: did the last campaign increase sales? does price affect switching?
  \item \textbf{Predict}: which bars will order more next quarter?
\end{itemize}

The second question is the sharp one, and it previews the whole second half: a brewer has sales data, not consumption data, and has it for its own products, not the competition's. Most of the interesting questions need data the company does not have, which is the found-data problem arriving from the commercial side.}

% ----------------------------------------------------
\begin{frame}
\frametitle{What is a research question?}
\centering

\begin{itemize}
  \item<1-> Any question we can answer
  \item<1-> Sometimes we say that we derive an RQ from a topic, and a theory from an RQ
  \begin{itemize}
    \item \textbf{Topic $>$ RQ $>$ Theory}
  \end{itemize}
  \item<2-> In reality, an RQ can be thought of as an operationalization of an argument
  \begin{itemize}
    \item \textbf{Previous evidence $>$ Argument $>$ RQ $>$ Hypotheses}
  \end{itemize}
  \item<2-> \textbf{Even though} this `argument' can be something anecdotical that we later develop into a proper, abstract theoretical argument
  \begin{itemize}
    \item<3-> And it would actually look into something like this:
    \item<3-> Previous $>$ `Anecdotal argument' $>$ RQ $>$ (Proper) Theory $>$ Hs
  \end{itemize}
\end{itemize}

\end{frame}
% ----------------------------------------------------
\note{The most useful twenty minutes in the course for students who have never written a design. Build the ladder slowly, one overlay at a time.

The first version is the one they will have been taught. The second is closer to how research actually works. The third is the honest one, and it is the one that gives them permission to start from a hunch, which is where everyone actually starts.

Callback to last week: the hotel and Spotify examples entered this ladder in the middle, at the data, and had to work backwards.}

% ----------------------------------------------------
% SWAPPABLE: current-event example. Keep the structure
% (topic > anecdotal argument > main RQ > RQ1/RQ2/RQ3 > theory);
% swap the topic for whatever is salient. The three sub-questions
% must split into descriptive / distributional / causal.
% ----------------------------------------------------
\begin{frame}
\frametitle{What is a research question? Example}
\centering

\begin{itemize}
  \item<1-> Topic/idea: Charlie Kirk and the rise of political violence?
  \item<2-> `Anecdotal argument'/intuition: online radicalization and young men more likely to commit individual acts of violence
  \item<3-> `Main' RQ: Does online media cause an increase in lone-wolf political violence?
  \begin{itemize}
    \item RQ1: Has political violence increased recently?
    \item RQ2: Political violence now more likely by individual young men?
    \item RQ3: Is online radicalization the cause of this?
  \end{itemize}
  \item<4-> Theoretical argument, hypotheses, expectations, etc
\end{itemize}

\end{frame}
% ----------------------------------------------------
\note{The pedagogical point is the three-way split, not the topic. Notice out loud what the sub-questions are:

\begin{itemize}
  \item RQ1 is \textbf{descriptive}: has it increased?
  \item RQ2 is \textbf{descriptive/distributional}: increased among whom?
  \item RQ3 is \textbf{causal}: is X the cause?
\end{itemize}

That is last week's triad reappearing inside a single project, and it is worth naming. Most students will present RQ3 as their question and then answer RQ1, which is the mismatch we diagnosed last week.

Ask which of the three is hardest, and why. The answer they usually give is RQ3; the honest answer is RQ1, and week 3 is about why.}

% ----------------------------------------------------
\begin{frame}
\frametitle{Good RQs, in brief}
\centering

\begin{itemize}
  \item[1.] Empirically \textbf{answerable}
  \begin{itemize}
    \item i.e. you can answer it with data
  \end{itemize}
  \item[2.] Theoretically \textbf{relevant}
  \begin{itemize}
    \item i.e. it helps you learn something about your theory/argument
  \end{itemize}
\end{itemize}

\end{frame}
% ----------------------------------------------------
\note{Two criteria, and both have to hold at once. Answerable but irrelevant is trivia; relevant but unanswerable is an essay, not a design.}

% ----------------------------------------------------
\begin{frame}
\frametitle{Good RQs, more in detail}
\centering

\begin{itemize}
  \item[1.] Consider potential results of the analyses
  \begin{itemize}
    \item if you found X, does that answer the question? causally?
    \item example: are kids who play videogames often more aggresive?
    \item does that inform a theory on the aggresiveness effect of VGs?
    \item[] (you can even try to do a better RQ without causal ID)
  \end{itemize}
  \item<2->[2.] Is it feasible?
  \begin{itemize}
    \item do you have the data? is it possible to do it? (e.g. re-offenders)
    \item also: is there any design or strategy to answer it?
  \end{itemize}
  \item<3->[3.] Keeping it simple and narrow
  \begin{itemize}
    \item what are the causes of economic underdevelopment? vs. does exposure to natural disasters hinders economic development?
  \end{itemize}
\end{itemize}

\end{frame}
% ----------------------------------------------------
\note{The first test is the most useful and the least practised: \textbf{imagine the result before you have it}. If every possible answer is uninteresting, the question is wrong, and you can know that today rather than in March.

Point 2 is where the second half of today lands: ``do you have the data'' is not a yes/no. It is ``does it exist, can you get it, and does it mean what you need it to mean''.}

% ----------------------------------------------------
\begin{frame}
\frametitle{Good RQs, more in detail}
\centering

\begin{itemize}
  \item[] \textbf{This:}
  \begin{itemize}
    \item Has political violence increased recently?
    \item Political violence now more likely by individual young men?
    \item Is online radicalization the cause of this?
  \end{itemize}
  \item[] \textbf{vs this:}
  \begin{itemize}
    \item What's the link between online media and political violence?
    \item[] {\color{asher}{\small (too broad)}}
    \item Is online media the cause of the increase of political violence among young men? {\color{asher}{\small (many questions in one)}}
    \item Why does online media drive political violence? {\color{asher}{\small (??)}}
  \end{itemize}
\end{itemize}

\end{frame}
% ----------------------------------------------------
\note{The three-way contrast is the best twenty seconds on the slide. Take them one at a time and make the room name the defect before you do.

The third one is worth dwelling on: ``why does X drive Y'' presupposes that X drives Y. A question that assumes its own answer cannot be tested, and students write these constantly.}

% ----------------------------------------------------
\begin{frame}
\frametitle{Example on generating RQs}
\centering

\begin{itemize}
  \item<1-> Couple things to remember:
  \begin{itemize}
    \item RQs are often the link between theory and empirics
    \item So they already suggest which variation to look at
  \end{itemize}
  \item[]
  \item<2-> Imagine you have the following argument:
  \item[]<2-> How good students do at school depends more on the peers they are surrounded by than on the quality of the teaching they receive
  \item[]
  \item<2-> Which RQ could let us test this?
\end{itemize}

\end{frame}
% ----------------------------------------------------
\note{Run it live, in pairs, three minutes, then take four or five answers. With forty people an open question dies; a pair turn does not.

What to listen for: most first attempts compare good schools with bad schools, which confounds peers and teaching completely. The useful move is to find variation that separates the two, which is a preview of weeks 4 and 5.

A good answer names the comparison, not just the topic.}

% ----------------------------------------------------
\begin{frame}
\frametitle{Stories, RQs, and theories}
\centering

\begin{itemize}
  \item There're no exclusive definitions of `stories' and theories
  \item It's just about getting to a sufficient level of abstraction
  \item Often, you start with a story or example, and then you move up the ladder for both theory and RQs until you get to a general theory tested with a RQ
\end{itemize}

\end{frame}
% ----------------------------------------------------
\note{The abstraction ladder. Students tend to think theory is something you either have or don't; it is something you climb towards from a story.}

% ----------------------------------------------------
\begin{frame}
\frametitle{Generating theories}
\centering

\begin{itemize}
  \item<1-> No recipe for this, everyone generates theories \textit{all the time}
  \item Usually it refers to an analytical argument that explains something
  \begin{itemize}
    \item It could also be a descriptive or predictive theory, but even in those cases there's probably an explanation underneath
  \end{itemize}
  \item Developed inductively, from descriptive data to general explanations
  \item My advice: if you can't tell a story out of the theory, you're not there yet {\small (i.e. need to be able to travel from/to abstraction)}
  \item<2->[]
  \item<2-> \textbf{Q:} How to identify a \textbf{good theory}?
\end{itemize}

\end{frame}
% ----------------------------------------------------
\note{The advice in the fourth bullet is the practical one and worth repeating: \textbf{if you cannot tell a story out of the theory, you are not there yet.} It is also a good diagnostic to apply in the workshop.

Take the closing question as a real question before turning the slide.}

% ----------------------------------------------------
\begin{frame}
\frametitle{Evaluating theories}
\centering

\begin{itemize}
  \item[1.]<1-> Simple
  \item[2.]<2-> Internally coherent and able to explain variation
  \item[3.]<3-> Testable
\end{itemize}

\end{frame}
% ----------------------------------------------------
\note{Three criteria, and the third is the one that does the work in this course. A theory that cannot fail cannot be tested, and a design that cannot fail is not a design.}

% ----------------------------------------------------
\begin{frame}
\frametitle{Example (of the whole process)}
\centering


\includegraphics[width = 0.45\textwidth]{../img/umu_erasmus}

\end{frame}
% ----------------------------------------------------
\note{A piece of descriptive evidence: Erasmus students do better in the labour market. Let them look at it and ask what they would conclude.}

% ----------------------------------------------------
\begin{frame}
\frametitle{Example}
\centering

\begin{itemize}
  \item That's some descriptive evidence that could inspire an anecdote
  \item The \textbf{anecdotal argument} (think of a story)
  \begin{itemize}
    \item<2-> My friend John who went on Erasmus has more money than my other friend who couldn't go, and also, John managed to get a job because his father is partner at a local firm
  \end{itemize}
  \item The \textbf{research question}
  \begin{itemize}
    \item<2-> Is there a causal effect of Erasmus on labor market early success? Is the effect mediated by household income?
  \end{itemize}
  \item The `proper' \textbf{theory}
  \begin{itemize}
    \item<2-> Going on Erasmus does not have any causal effect on getting a first job, the relationship is explained by the confounding effect of income
    \item<2-> Or: Positive effect among high-income students because they have access to informal networks where this experience is valued
  \end{itemize}
  \item The \textbf{hypotheses}?
\end{itemize}

\end{frame}
% ----------------------------------------------------
\note{The whole ladder in one example, and the two candidate theories are the point: the same descriptive fact supports a ``no effect, it is all income'' story and a ``real effect, but only for some'' story.

Those are different theories with different observable implications, which is exactly what makes them testable. Ask what you would have to see to tell them apart.}

% ----------------------------------------------------
\begin{frame}
\frametitle{Hypotheses}
\centering

\begin{itemize}
  \item Hypothesis is just a very formal term for empirical expectations
  \begin{itemize}
    \item Which essentially means being able to say what you expect to see given a theory
    \item[]<2-> And ideally, knowing what you need to see in order to discredit the theory
  \end{itemize}
  \item[]<3->
  \item<3-> Imagine that I have the theory that my knee hurts when I do sports on cold days
  \begin{itemize}
    \item Simplifying it, we have a 2x2 situation:
  \end{itemize}
\end{itemize}

\end{frame}
% ----------------------------------------------------
\note{The second sub-bullet is the important half and the one students skip: what would have to be true for you to abandon the theory?

Anyone who cannot answer that does not have a hypothesis, they have a hope. This comes back in week 5 with pre-registration.}

% ----------------------------------------------------
\begin{frame}
\frametitle{Hypotheses: what would you expect to observe?}
\centering

\begin{tabular}{|l|c|c|}
\toprule
            & \textit{Cold day} & \textit{Hot day} \\
\midrule
\textit{Run}         &   \only<2>{Pain}   & \only<2>{Not pain} \\
\midrule
\textit{Didn't run}  & \only<2>{Not pain} & \only<2>{Not pain} \\
\bottomrule
\end{tabular}

\end{frame}
% ----------------------------------------------------
\note{Make them fill in the four cells before you reveal. Pairs, one minute.

The theory is ``cold \textit{and} running'', so only one cell should show pain. Getting them to see that a theory implies a whole pattern, not a single observation, is the point.}

% ----------------------------------------------------
\begin{frame}
\frametitle{What if you observe this? \textbf{New theory? Test?}}
\centering

\begin{tabular}{|l|c|c|}
\toprule
            & \textit{Cold day} & \textit{Hot day} \\
\midrule
\textit{Run}         & Pain & Not pain \\
\midrule
\textit{Didn't run}  & Pain & Pain \\
\bottomrule
\end{tabular}

\end{frame}
% ----------------------------------------------------
\note{Now the data disagrees with the theory, and in an informative way: pain on cold days regardless of running, and no pain when running on hot days. Running looks irrelevant; cold does not.

Two things worth drawing out: the theory was wrong in a \textit{specific} way, which tells you what to think next; and you could only learn that because the theory made a prediction for all four cells. A vaguer theory would have survived.}

% ====================================================
\section{Theory, mechanisms, levels}
% ====================================================

% ----------------------------------------------------
\begin{frame}
\frametitle{Computational methods and theory}
\centering

\begin{itemize}
  \item Limitations of \textit{data mining}
  \item Focus on the what rather that on the why
  \item Problems with machine learning
  \begin{itemize}
    \item Example of predicting ice cream sales
  \end{itemize}
\end{itemize}

\end{frame}
% ----------------------------------------------------
\note{The obvious objection to everything so far: why start with questions and arguments at all, when you could just let the data speak?

Answer at design level, not as a complaint about machine learning. A model that predicts ice cream sales from drowning rates works fine and tells you nothing you can act on. Prediction without a why is useful (last week's Blumenstock example) but it answers a different question, and confusing the two is the mismatch from week 1.

Relevant here because half this cohort arrived wanting to do exactly the thing this slide qualifies.}

% ----------------------------------------------------
\begin{frame}
\frametitle{Why bother about theory?}
\centering

\includegraphics[width = \textwidth]{../img/predicting_conflict}

\end{frame}
% ----------------------------------------------------
\note{Predicting conflict: a well-resourced effort, plenty of data, and performance that is much worse than the enthusiasm suggested.

The design lesson is not ``prediction is bad''. It is that rare events, changing processes and outcomes that get intervened on are hard to predict, and no amount of data fixes a problem that lives in the data-generating process. Week 5 closes on the same theme with Fragile Families.}

% ----------------------------------------------------
\begin{frame}
\frametitle{Mechanisms briefly}
\centering

\begin{itemize}
  \item A \textbf{mechanism} is basically the \textbf{how} (or why) of a relationship
  \item e.g., we know that the flu gives us fever
  \begin{itemize}
    \item flu $>$ fever
  \end{itemize}
  \item<2-> What's the mechanism?
  \begin{itemize}
    \item flu $>$ \textit{immune system detects infection} $>$ $\Delta$ \textit{body temp} $>$ fever
  \end{itemize}
  \item<3-> Good thing about mechanisms is that we can try to test them
\end{itemize}

\end{frame}
% ----------------------------------------------------
\note{Mechanisms matter here for a practical reason: they generate sub-questions, and sub-questions are usually more answerable than the main one.

If you cannot get at ``does X cause Y'', you can often get at one step of the chain, and that is a real contribution rather than a failure.}

% ----------------------------------------------------
\begin{frame}
\frametitle{Testing mechanisms ($\approx$ sub-research questions)}
\centering

\begin{itemize}
  \item Let's go back to the Erasmus example
  \item If our theory is that the effect of going on Eramus is higher for high-income students due to their access to social networks,
  \item[] \textbf{what's the mechanism?}
  \item And how could we try to \textbf{test it}?
  \begin{itemize}
    \item (Think about the sub-RQs)
  \end{itemize}
\end{itemize}

\end{frame}
% ----------------------------------------------------
\note{Pairs, three minutes.

The mechanism is informal networks. Sub-questions that follow: do high-income Erasmus students end up in different sectors? do they report using personal contacts to find work more often? does the effect disappear where hiring is anonymised?

That last one is a design idea, not just a question, and it is worth pointing out that they just did the thing week 5 is about.}

% ----------------------------------------------------
\begin{frame}
\frametitle{Complexity of the social world and micro/macro}
\centering

\includegraphics[width = 0.8\textwidth]{../img/leipzig1989}

\begin{itemize}
  \item Why do mass protests emerge?
\end{itemize}

\end{frame}
% ----------------------------------------------------
\note{Leipzig, 1989. Moved here from the measurement material, because levels of explanation belong with theory-building rather than with measurement.

Ask the question and let it sit. Answers will mix levels freely: grievances, opportunity, someone being first, the state hesitating. That mixing is the material for the next slide.}

% ----------------------------------------------------
\begin{frame}
\frametitle{Levels of explanation}
\centering

\begin{itemize}
\item[] Can you think of...?
\item Macro-level mechanisms
\item Micro-level mechanisms
\item[]
\item<2-> What's the point of macro-level explanations, actually?
\end{itemize}

\end{frame}
% ----------------------------------------------------
\note{Macro: economic decline, regime weakness, diffusion from Poland and Hungary. Micro: an individual deciding to turn up, given what they expect others to do.

The closing question is the good one. Macro explanations are not micro explanations in disguise; they can be about aggregate patterns that no individual produces on their own. But a macro claim with no plausible micro story underneath is usually a correlation with ambition.

Connect it to the unit of analysis from week 1: choosing a level is choosing a question.}

% ----------------------------------------------------
\begin{frame}
\frametitle{Recap}
\centering

\begin{itemize}
  \item Questions on theory, RQ, or mechanisms?
\end{itemize}

\end{frame}
% ----------------------------------------------------
\note{Genuine pause before the break. The second half is a different topic and it is worth arriving there clean.

Break here, ten minutes.}

% ====================================================
\section{Where does the data come from?}
% ====================================================

% ----------------------------------------------------
\begin{frame}
\frametitle{Last week, in one line}
\centering

\vspace{15pt}
\begin{columns}[T]
\begin{column}{0.47\textwidth}
\centering \textbf{Designed}
\vspace{4pt}
\hrule
\vspace{6pt}
\footnotesize
\begin{flushleft}
answers a question \textbf{you} had
\end{flushleft}
\end{column}
\begin{column}{0.47\textwidth}
\centering \textbf{Found}
\vspace{4pt}
\hrule
\vspace{6pt}
\footnotesize
\begin{flushleft}
answers a question \textbf{someone else} had.\\Or no question at all
\end{flushleft}
\end{column}
\end{columns}

\vspace{25pt}
\uncover<2->{\centering Today: what follows from being on the right-hand side.}

\end{frame}
% ----------------------------------------------------
\note{Thirty seconds of recall, then straight on. Everything in this half is a consequence of the right-hand column.}

% ----------------------------------------------------
\begin{frame}[plain]
  \begin{tikzpicture}[remember picture, overlay]
    \node[at = (current page.center), yshift = 0cm] (cover) {%
      \includegraphics[keepaspectratio, width=\paperwidth, height=\paperheight]{../img/salganik_readymade}};
  \end{tikzpicture}
\end{frame}
% ----------------------------------------------------
\note{Salganik's framing, and the best one available: \textbf{readymade} and \textbf{custommade}.

Duchamp's \textit{Fountain} is a urinal he did not make, seen and repurposed. Michelangelo's \textit{David} was made for the purpose. Both are art; they involve completely different skills.

Found data is a readymade. The creative work is not in producing it, it is in recognising what an object built for someone else's purpose can and cannot be used for. That is a real skill and it is the one this half of the session is about.}

% ----------------------------------------------------
\begin{frame}
\frametitle{Designed data: someone decided what to record}
\centering

\begin{minipage}{.48\textwidth}\centering
\includegraphics[keepaspectratio, width=\textwidth, height=0.55\paperheight]{../img/afrobaro}
\end{minipage}\hfill
\begin{minipage}{.48\textwidth}\centering
\includegraphics[keepaspectratio, width=\textwidth, height=0.55\paperheight]{../img/ucdp_info}
\end{minipage}

\vspace{6pt}
{\footnotesize A survey, and a hand-coded conflict dataset. Someone wrote the questions, or the coding rules.}

\end{frame}
% ----------------------------------------------------
\note{Two examples of custommade data: a cross-national survey, and an event dataset built by researchers reading sources and applying coding rules.

The property that matters: there is a \textbf{codebook}, and it exists because someone had to decide in advance what counted. Those decisions are documented, arguable, and revisable. You can go and read why a given event is or is not in the data.

Expensive, slow, and small. That is the trade.}

% ----------------------------------------------------
\begin{frame}
\frametitle{Found data: nobody decided anything, for you}
\centering

\begin{minipage}{.48\textwidth}\centering
\includegraphics[keepaspectratio, width=\textwidth, height=0.55\paperheight]{../img/nature_creditcard}
\end{minipage}\hfill
\begin{minipage}{.48\textwidth}\centering
\includegraphics[keepaspectratio, width=\textwidth, height=0.55\paperheight]{../img/twitter_protests}
\end{minipage}

\vspace{6pt}
{\footnotesize Credit card transactions and social media posts, repurposed as research data.}

\end{frame}
% ----------------------------------------------------
\note{Two readymades. Credit card transactions exist so that banks can settle payments; posts exist because a platform wanted engagement.

Neither has a codebook, because nobody was answering a research question when they were produced. There is no document explaining what counts as a purchase or a protest post, because nobody needed one.

Big, fast, cheap, and constantly available. That is the other side of the trade, and it is why everyone wants to use them.}

% ----------------------------------------------------
\begin{frame}
\frametitle{The asymmetry that matters}
\centering

\vspace{10pt}
\begin{itemize}[<+->]
  \item With \textbf{designed} data, you know what the numbers were meant to mean, because someone wrote it down
  \vspace{10pt}
  \item With \textbf{found} data, you have to \textbf{reconstruct} it, from the purpose of whoever built the system
  \vspace{12pt}
  \item So: \colorbox{yellow}{who made this, and what were they trying to do?}
\end{itemize}

\end{frame}
% ----------------------------------------------------
\note{The sentence to land, and the one that generalises to their whole MA.

Facebook did not build a friendship graph so you could measure social capital. Spotify did not tag genres so you could study cultural convergence. The purpose is still in the data, and if you do not reconstruct it you will mistake a fact about the system for a fact about the world.

This is exactly the Spotify question from last week: is this a fact about music, or a fact about Spotify?

It is also Memo 2, so say it slowly.}

% ====================================================
\section{Ten things about big data sources}
% ====================================================

% ----------------------------------------------------
\begin{frame}
\frametitle{Salganik's ten characteristics}
\centering

\vspace{8pt}
\begin{columns}[T]
\begin{column}{0.45\textwidth}
\centering \textbf{Why we want them}
\vspace{4pt}
\hrule
\vspace{8pt}
\footnotesize
\begin{flushleft}
Big\\[4pt]
Always-on\\[4pt]
Non-reactive
\end{flushleft}
\end{column}
\begin{column}{0.52\textwidth}
\centering \textbf{Why designs break}
\vspace{4pt}
\hrule
\vspace{8pt}
\footnotesize
\begin{flushleft}
Incomplete\\[3pt]
Inaccessible\\[3pt]
Non-representative\\[3pt]
Drifting\\[3pt]
Algorithmically confounded\\[3pt]
Dirty\\[3pt]
Sensitive
\end{flushleft}
\end{column}
\end{columns}

\vspace{12pt}
{\footnotesize \asher{Salganik, \textit{Bit by Bit}, ch. 2. Free online.}}

\end{frame}
% ----------------------------------------------------
\note{A checklist they will use for the rest of the MA, and there is not a single statistical idea in it.

Three on the left, seven on the right: that ratio is the honest summary of working with found data. Do not present the right-hand column as reasons not to use it. Every one of them is manageable \textit{if you know it is there before you design}, and fatal if you discover it in March.

We take four of the seven properly today: non-representative, drifting, algorithmically confounded, and inaccessible. Dirty gets a mention. Incomplete and sensitive point forward to weeks 3 and 5.}

% ----------------------------------------------------
\begin{frame}
\frametitle{Big, always-on, non-reactive}
\centering

\begin{itemize}[<+->]
  \item \textbf{Big}: rare events become studiable, and small differences become visible
  \begin{itemize}
    \item but big does not fix bias. It makes you \textit{confident} about the wrong number
  \end{itemize}
  \vspace{6pt}
  \item \textbf{Always-on}: you can study things that were not anticipated, and study them \textit{before and after}
  \begin{itemize}
    \item unexpected events get a baseline for free
  \end{itemize}
  \vspace{6pt}
  \item \textbf{Non-reactive}: people are not answering your questions, so they are not performing for you
  \begin{itemize}
    \item no social desirability bias. But no informed consent either
  \end{itemize}
\end{itemize}

\end{frame}
% ----------------------------------------------------
\note{The three genuine attractions, each with the qualification attached, because the qualification is where the design work is.

\textbf{Big}: the sub-bullet is the one to stress. More data narrows the confidence interval around a biased estimate. It buys precision, not validity. Callback to Challenger: no number of failures would have helped.

\textbf{Always-on}: the strongest and least-appreciated one. If you had wanted to study behaviour before and after some sudden event, no survey would have been in the field beforehand. Platform data always is.

\textbf{Non-reactive}: cuts both ways, and the second half is the seed of the ethics discussion in week 5.}

% ----------------------------------------------------
\begin{frame}
\frametitle{Non-representative}
\centering

\vspace{10pt}
\begin{itemize}[<+->]
  \item Who is on the platform is not who is in the population
  \item Who \textit{posts} is not who is on the platform
  \vspace{12pt}
  \item So: is a non-representative sample useless?
\end{itemize}

\end{frame}
% ----------------------------------------------------
\note{Two layers of selection, and students routinely notice the first and miss the second. The people who post about politics on any platform are a tiny, unusual subset of the people on it.

Then ask the closing question and take answers before turning the page. Most rooms split into ``yes, useless'' and ``no, big enough to be fine''. Both are wrong, which is what makes the next example worth showing.}

% ----------------------------------------------------
\begin{frame}
\frametitle{The Xbox poll}
\centering

\vspace{5pt}
\begin{itemize}[<+->]
  \item 2012 US election. Opt-in poll served to \textbf{Xbox} users
  \item $\sim$350,000 respondents. Around 93\% male, overwhelmingly young
  \vspace{8pt}
  \item Raw numbers: a large Romney lead. Wrong, and obviously so
  \vspace{8pt}
  \item Adjusted using what was \textit{known} about who was missing:\\ tracked the campaign closely and called the outcome
\end{itemize}

\vspace{10pt}
\uncover<5->{{\footnotesize \asher{Wang, Rothschild, Goel \& Gelman (2015), \textit{Int.\ J.\ Forecasting} 31(3).}}}

% TODO: source the demographic-skew bar chart and the forecast-vs-outcome plot
\end{frame}
% ----------------------------------------------------
\note{An absurdly unrepresentative sample producing a good forecast.

\textbf{The lesson is not ``non-representative is fine''.} It is that representativeness is a property of the \textbf{design}, not of the sample. What made this work was knowing precisely \textit{how} the sample differed, and having something external to adjust against.

{\footnotesize Two consequences, cutting in opposite directions, and students usually hold only one:
\begin{itemize}
  \item ``My Twitter data isn't random so it's useless'': not necessarily, if you can characterise the skew.
  \item ``I have ten million observations so it doesn't matter'': it matters more, and you will be very precisely wrong.
\end{itemize}}

The adjustment is statistics and we are not doing it. Our question is: \textit{do you know who is missing, and do you have anything to anchor to?} If not, the design fails at any sample size.}

% ----------------------------------------------------
\begin{frame}
\frametitle{Drifting}
\centering

\vspace{15pt}
\begin{itemize}[<+->]
  \item The world changes
  \item The people on the platform change
  \item \textbf{The platform changes}
\end{itemize}

\vspace{20pt}
\uncover<4->{\centering A measure can stop meaning what it meant, \colorbox{yellow}{without anything looking wrong}.}

\end{frame}
% ----------------------------------------------------
\note{Three sources of drift, and the third is the one that is specific to found data and invisible from inside the dataset.

The closing line is the whole problem. A drifting measure does not throw an error. The numbers keep arriving, the series keeps updating, and the trend you are looking at is partly a trend in the instrument.}

% ----------------------------------------------------
\begin{frame}
\frametitle{Google Flu Trends}
\centering

\vspace{5pt}
\begin{itemize}[<+->]
  \item 2008: predict flu prevalence from search queries. Faster than the CDC, and it worked
  \vspace{6pt}
  \item 2009: missed the H1N1 pandemic, a flu outbreak that was not seasonal
  \vspace{6pt}
  \item 2011--13: over-predicted almost every week. In the 2012--13 season, more than \textbf{double} the CDC figure
\end{itemize}

\vspace{12pt}
\uncover<4->{{\footnotesize \asher{Lazer, Kennedy, King \& Vespignani (2014), \textit{Science} 343(6176). This week's paper.}}}

% TODO: source the CDC-vs-GFT over-prediction time series figure
\end{frame}
% ----------------------------------------------------
\note{Do not give away the diagnosis yet. The paper discussion is later tonight and this is its setup.

The sequence matters: it worked, then it failed in a specific way, then it failed persistently. A system that had never worked would not be interesting.

If someone asks why, put it back to the room and hold it for the discussion.}

% ----------------------------------------------------
\begin{frame}
\frametitle{What went wrong?}
\centering

\vspace{10pt}
\begin{itemize}[<+->]
  \item \textbf{Big data hubris}: the assumption that volume substitutes for measurement
  \begin{itemize}
    \item it was fitted to correlate with CDC data, not to measure flu
  \end{itemize}
  \vspace{10pt}
  \item \textbf{Algorithm dynamics}: Google kept changing the search engine
  \begin{itemize}
    \item suggested searches, autocomplete, health panels
    \item people searched differently because the \textit{product} changed
  \end{itemize}
\end{itemize}

\vspace{10pt}
\uncover<3->{\centering The instrument was being modified, \colorbox{yellow}{by someone who was not doing research}.}

\end{frame}
% ----------------------------------------------------
\note{Lazer et al.'s two named causes.

\textbf{Big data hubris} is the one with the wider reach: the model was selected to correlate with CDC reports. Some of the winning search terms were seasonal rather than flu-related, which is fine until a flu season arrives out of season, as in 2009.

\textbf{Algorithm dynamics} is the CSS-specific one, and it is the deeper problem. Google was improving its product, as it should. Every improvement changed the behaviour being measured. Nobody at Google was doing anything wrong, and the measure broke anyway.

The closing line is the transferable one: with found data, \textbf{your instrument is owned and maintained by someone whose interests are not yours}, and they will change it without telling you.

Direct line to Spotify retagging genres last week. Say so.}

% ----------------------------------------------------
\begin{frame}
\frametitle{Algorithmically confounded}
\centering

\vspace{10pt}
\begin{itemize}[<+->]
  \item Platforms do not just record behaviour. They \textbf{shape} it
  \begin{itemize}
    \item ranking, recommendation, defaults, what is easy and what is hidden
  \end{itemize}
  \vspace{10pt}
  \item So an observed pattern may be a property of \textbf{the system}, not of the people
  \vspace{10pt}
  \item The platform is in your causal diagram \hfill \asher{(week 4)}
\end{itemize}

\end{frame}
% ----------------------------------------------------
\note{Distinct from drift, and students merge the two. Drift is the instrument changing over time. Algorithmic confounding is the instrument \textit{causing} the thing you are measuring, even when nothing changes.

Concrete: friendship networks on a platform that recommends friends partly reflect the recommender. Popularity measured where popular things are shown more will show more popularity. You cannot separate the behaviour from the system that elicited it by looking harder at the data.

Forward-reference DAGs in week 4 explicitly. When they draw a diagram for a platform study, the platform is a node.

If you have time: this is the strongest argument for why an \textit{experiment} on a platform (Guess et al., week 4) is worth so much more than observation on one.}

% ----------------------------------------------------
\begin{frame}
\frametitle{Inaccessible, and the post-API age}
\centering

\vspace{8pt}
\begin{itemize}[<+->]
  \item The data belongs to someone, and they can revoke access
  \begin{itemize}
    \item Twitter/X academic access, closed 2023
    \item CrowdTangle, shut down 2024
  \end{itemize}
  \vspace{8pt}
  \item A design that depends on access you do not control is a design that can \textbf{stop working}
  \vspace{8pt}
  \item And nobody can check your results afterwards
\end{itemize}

\vspace{8pt}
\uncover<4->{{\footnotesize \asher{Freelon (2018), Computational Research in the Post-API Age, \textit{Political Communication} 35(4).}}}

\end{frame}
% ----------------------------------------------------
\note{Practical, and several of them will propose scraping something for the final essay, so say it plainly.

Two distinct problems. \textbf{Feasibility}: can you still get it in March? \textbf{Reproducibility}: if the data cannot be obtained again, nobody can check the finding, including you.

The design question for Memo 2 and for the essay: \textit{can you actually get this, and could someone else get it again?} A project that fails this is better abandoned in week 2 than in April.

Do not moralise. The honest position is that this is a real constraint on what computational social science can currently be, not a failure of nerve on the students' part.}

% ----------------------------------------------------
\begin{frame}
\frametitle{Dirty, and sensitive}
\centering

\vspace{10pt}
\begin{itemize}[<+->]
  \item \textbf{Dirty}: spam, bots, duplicates, tests, automated accounts
  \begin{itemize}
    \item the cleaning decisions are design decisions, and they are rarely reported
  \end{itemize}
  \vspace{12pt}
  \item \textbf{Sensitive}: much of it is about identifiable people who never agreed to anything
  \begin{itemize}
    \item ``it's public'' is not the same as ``they consented''
    \item aggregation and linkage can re-identify \hfill \asher{(week 5)}
  \end{itemize}
\end{itemize}

\end{frame}
% ----------------------------------------------------
\note{Two minutes, both flagged rather than covered.

\textbf{Dirty}: the point is not that data needs cleaning, which they know. It is that ``we removed bots'' is a choice with consequences that almost never appears in a paper, and different reasonable choices give different answers. That connects to week 5 on researcher degrees of freedom.

\textbf{Sensitive}: flag it now, deliberately, because essay proposals get written in week 4 and the ethics session is week 5. If someone is planning to scrape and link personal data, they should meet the objection before they have invested in the design, not after.

The credit-card study earlier in the session is a good callback: entirely legitimate research, and also a dataset that could identify almost anyone in it.}

% ----------------------------------------------------
\begin{frame}
\frametitle{Incomplete}
\centering

\vspace{10pt}
\begin{itemize}[<+->]
  \item The variable you need is almost never the variable that was recorded
  \begin{itemize}
    \item platforms record what is useful to them, not what is useful to you
  \end{itemize}
  \vspace{10pt}
  \item No demographics. No income. No control variables
  \vspace{10pt}
  \item You can sometimes \textbf{infer} what is missing. That is a model, and models can be wrong
\end{itemize}

\end{frame}
% ----------------------------------------------------
\note{The seventh characteristic, and the one that bites earliest in a student project.

Concretely: a platform dataset will tell you what an account posted and when, and almost never who the person is, how much they earn, or where they live. The variables you would want as controls are exactly the ones nobody had a commercial reason to store.

The third bullet is the bridge to next week. Inferring gender or ideology or location from behaviour is common and sometimes done well, but it turns one measurement problem into two: now you need to know how good the inference is, and the error will not be random. That is the classifier-output-as-a-variable problem, and it is session 3.}

% ====================================================
\section{So where do you actually get data?}
% ====================================================

% ----------------------------------------------------
\begin{frame}
\frametitle{A practical interlude}
\centering

\vspace{20pt}
{\large You have a question. Now what?}

\vspace{25pt}
\begin{itemize}
  \item<2-> Most projects in this programme use data somebody has already published
  \item<2-> Knowing what exists, and what it looks like when it arrives, is a real skill
\end{itemize}

\end{frame}
% ----------------------------------------------------
\note{Fifteen minutes, and unglamorous on purpose. It is also the most immediately useful thing tonight, and it is exactly what Memo 2 asks for.

Students routinely propose designs that require data nobody has ever collected, and equally routinely miss that the thing they need has been sitting on a public portal for fifteen years. Both are failures of the same skill.}

% ----------------------------------------------------
\begin{frame}
\frametitle{Places things actually live}
\centering

\begin{minipage}{.48\textwidth}\centering
\includegraphics[keepaspectratio, width=\textwidth, height=0.5\paperheight]{../img/wbank}
\end{minipage}\hfill
\begin{minipage}{.48\textwidth}\centering
\includegraphics[keepaspectratio, width=\textwidth, height=0.5\paperheight]{../img/eurostat}
\end{minipage}

\vspace{8pt}
{\footnotesize Also: INE, Eurobarometer, ESS, V-Dem, UCDP, national statistical offices, municipal open data portals, replication archives (Dataverse, OSF).}

\end{frame}
% ----------------------------------------------------
\note{Do not read the list. Point out the categories and move:

\begin{itemize}
  \item \textbf{International organisations} (World Bank, Eurostat, OECD): comparable, aggregated, usually country-year.
  \item \textbf{National statistical offices} (INE): far more granular, often only in Spanish, often municipal.
  \item \textbf{Academic projects} (V-Dem, UCDP, ESS, Afrobarometer): built for research, with codebooks.
  \item \textbf{Replication archives}: the most underused. Every paper in a decent journal now deposits its data, which means the measure someone built for their paper is sitting there for you.
\end{itemize}

That last one is worth thirty seconds. If a published paper measured the thing they care about, they can usually just get it.}

% ----------------------------------------------------
\begin{frame}
\frametitle{What a codebook gives you}
\centering

\includegraphics[keepaspectratio, width=0.92\textwidth, height=0.72\paperheight]{../img/afrobaro_codebook}

\end{frame}
% ----------------------------------------------------
\note{Afrobarometer. This is what designed data looks like from the inside, and it is worth reading a couple of lines aloud.

Three things to point at:
\begin{itemize}
  \item The \textbf{exact question wording}. ``How much do you trust the police'' is a specific stimulus, not a concept. Next week is about the gap between the two.
  \item The \textbf{value labels}, including 9 for ``don't know'' and 98 for ``refused''. If you analyse this without reading the codebook, you will average in a 98 and produce nonsense. This happens constantly.
  \item The footnote: \textit{not asked in BFO, SWZ}. The variable silently does not exist for two countries, and nothing in the data file will tell you.
\end{itemize}

Found data has none of this. That is the whole point of the section.}

% ----------------------------------------------------
\begin{frame}
\frametitle{And what it looks like when it arrives}
\centering

\includegraphics[keepaspectratio, width=0.95\textwidth, height=0.62\paperheight]{../img/electoral_data_excel}

\vspace{6pt}
{\footnotesize Spanish municipal election results. What is the unit? What is missing?}

\end{frame}
% ----------------------------------------------------
\note{Spend two or three minutes here, as a live reading exercise. Ask both questions on the slide and take answers.

\begin{itemize}
  \item \textbf{Unit}: the municipality. But each row holds two elections, so it is really municipality-wide with election-specific columns. To analyse change you would reshape it.
  \item \textbf{Missing}: look at Orbaizeta, which has no PP vote share in 2019. Is that a zero, a coalition running under another name, or a genuine gap? \textbf{The file will not tell you.} A blank cell and a real zero look identical once you load it.
  \item The columns are shares, not counts, so you cannot recover turnout numbers or weight by size without merging population in.
\end{itemize}

The transferable lesson: \textbf{a dataset is not self-explanatory}, and every one of these questions has to be answered before the first analysis, not after.}

% ----------------------------------------------------
\begin{frame}
\frametitle{Before you commit to a data source}
\centering

\vspace{5pt}
\begin{itemize}[<+->]
  \item Who made it, and \textbf{why}?
  \item What is the \textbf{unit}, and what is the coverage in space and time?
  \item Is there a \textbf{codebook}? If not, what replaces it?
  \item Who or what is \textbf{missing}?
  \item Can you actually \textbf{get it}, and could someone else get it again?
  \item Has it \textbf{changed} over the period you care about?
\end{itemize}

\vspace{12pt}
\uncover<7->{{\footnotesize \asher{Six questions. They are also, more or less, Memo 2.}}}

\end{frame}
% ----------------------------------------------------
\note{The checklist, and the practical distillation of the whole evening. Every question maps back to something from the session: purpose to the readymade/custommade distinction, missing to non-representative and incomplete, getting it to post-API, changed to drift.

Tell them to run any dataset they are considering for the essay through these six. A source that fails question 5 is not a source, however good the idea.

This is the slide to put back up when you set Memo 2 in twenty minutes.}

% ====================================================
\section{Paper: The Parable of Google Flu}
% ====================================================

% ----------------------------------------------------
\begin{frame}
\frametitle{Why this paper}
\centering

\vspace{20pt}
\begin{itemize}
  \item Three pages, no equations
  \vspace{8pt}
  \item Measurement validity, drift, and algorithmic confounding \textbf{in one case}
  \vspace{8pt}
  \item And a well-resourced team getting it wrong in public
\end{itemize}

\end{frame}
% ----------------------------------------------------
\note{One slide on why it is on the list, per the structure we use every week: this is the design question the paper illustrates.

Then the sequence: I set it up, pairs work the prompts on screen, small groups redesign it, I close. About twenty-five minutes in total, so keep it moving.

Nothing was handed in and nothing is marked, so start by asking how many have read it and adjust honestly rather than pretending.}

% ----------------------------------------------------
\begin{frame}
\frametitle{In pairs}
\centering

\vspace{10pt}
\begin{itemize}
  \item[1.] What was GFT actually \textbf{measuring}?
  \vspace{10pt}
  \item[2.] Whose behaviour changed: the public's, or Google's?
  \vspace{10pt}
  \item[3.] Would \textbf{more data} have fixed it?
  \vspace{10pt}
  \item[4.] What would you have had to do in 2009 to see this coming?
\end{itemize}

\end{frame}
% ----------------------------------------------------
\note{\textbf{Twelve minutes total}, so this has to move. Prompts 1 and 2 together: three minutes in pairs, then three or four contributions. Then 3 and 4 the same way. Open-floor does not work at this size, but neither does letting a pair turn run long.

\begin{itemize}
  \item[1.] Not flu. Search behaviour correlated with CDC reports of flu. The gap between those is the whole paper.
  \item[2.] Both, and that is the point. Google changed the product; users changed how they searched in response.
  \item[3.] No. More of a mis-specified measure is more of a mis-specified measure. This is the ``big'' qualification from earlier.
  \item[4.] The strongest answers: validate against something external and keep validating; treat the search engine as a changing instrument; report the model rather than keeping it proprietary.
\end{itemize}}

% ----------------------------------------------------
\begin{frame}
\frametitle{Groups of four: redesign it}
\centering

\vspace{20pt}
{\large You have Google's search logs and CDC data, in 2009.}

\vspace{15pt}
{\large Build something that does not break the same way.}

\vspace{20pt}
{\footnotesize \asher{Seven minutes. One change, and what it buys you.}}

\end{frame}
% ----------------------------------------------------
\note{The step that matters most. Asking students to \textit{evaluate} a Science paper produces deference; asking what they would change produces engagement, and it is the same skill the workshop tests.

Good answers, roughly in order of sophistication: combine search with CDC data rather than replacing it; re-fit continuously instead of once; choose terms with a causal story rather than by correlation alone; hold out seasons rather than weeks; publish the model so failures are diagnosable.

Take three groups, not twelve. Seven minutes, and cut to five if the prompts ran over.}

% ----------------------------------------------------
\begin{frame}
\frametitle{The design lesson}
\centering

\vspace{25pt}
{\large Your instrument was built by someone else,\\for their own purposes,}

\vspace{10pt}
{\large and they will keep changing it.}

\vspace{25pt}
{\footnotesize \asher{Which is a measurement problem. That is next week.}}

\end{frame}
% ----------------------------------------------------
\note{Close on one sentence and hand to week 3.

Google Flu is usually told as a story about hubris. The more useful reading for this cohort is that it is a \textbf{measurement failure}: nobody ever established that the thing being counted was the thing they cared about, and no amount of data or engineering rescues that.

That is precisely question 3 of the course, which is next week.}

% ====================================================
\section{Before next week}
% ====================================================

% ----------------------------------------------------
\begin{frame}
\frametitle{Before next week}
\centering

\vspace{8pt}
\begin{itemize}
  \item \textbf{Read:} Müller-Crepon, Hunziker \& Cederman, \href{https://journals.sagepub.com/doi/10.1177/0022002720963674}{\textit{Roads to Rule, Roads to Rebel}}
  \begin{itemize}
    \item Longer than Google Flu. Skip the statistics, read the measurement
  \end{itemize}
  \vspace{12pt}
  \item<2-> \textbf{Memo 2}, in pairs, Aula Global by \textbf{Monday night}:
  \begin{itemize}
    \item<2-> Take a source of found data you could realistically get hold of
    \item<2-> Who produced it, and why? Who or what is missing from it?
    \item<2-> Name one question it could answer well, and one it could not
    \item<2-> Half a page
  \end{itemize}
\end{itemize}

\end{frame}
% ----------------------------------------------------
\note{Memo 2 is this whole half-session applied to something they choose. Push them towards data they could genuinely obtain, since the feasibility question from earlier is the point, and a memo about a dataset they will never get is a wasted exercise.

Good picks: a public API they have already used, an institutional dataset, municipal open data, something scraped from a site with sane terms. Bad picks: anything requiring platform access they do not have.

Remind them it is not about their final essay, and that pairs can change.

On the reading: be honest that it is a proper journal article rather than three pages, and tell them exactly what to read it for: the concept, and how it gets turned into a measure. They can skip the models entirely.}

% ----------------------------------------------------
\begin{frame}
\frametitle{}
\centering

\vspace{40pt}
{\Large Questions?}

\end{frame}
% ----------------------------------------------------
\note{}
